% Einheit 4 von 4 – Parallele Ebenen (bank/ebenen/e4.jsonl, zusammenbau v0.1)
%% TODO Zeitmarke Einheit 4: Marken-Zeile nicht lesbar
%% TODO Zweigzeile Teil 1: was hier gelernt wird (Satz in Schülersprache aus dem Einheitstitel) – Einheit 4
\einheitenkopf[e4]{Einheit 4 von 4 · Parallele Ebenen}
\zweigzeile{Abitur GK}
\setcounter{aufgabe}{16}

%% TODO Titel: Ich-kann-Satz fehlt in der Bank (Platzhalter: Kettenname) – Nr. 17
%% TODO Anweisung: ein Satz über der ersten Teilaufgabe fehlt in der Bank – Nr. 17
\begin{aufgabe}{Parallele Ebenen}
%% TODO \rechenplatz{n}: Schrittzahl des Grundfalls fehlt in der Bank (nur bei fester Schreibform, 2.3 b)
\begin{teile}
\teil $E\colon x + 2y - z = 3$ und $F\colon 2x + 4y - 2z = 6$. Wie liegen sie zueinander? \\ \kreuz{parallel} \\ \kreuz{identisch} \\ \kreuz{schneidend}
\teil Ebene $F$ parallel zu $E\colon 3x + y - 2z = 4$ durch $P(2 | 1 | 1)$?
\teil Ein Würfel mit der Kantenlänge 6 hat eine Ecke im Ursprung und liegt im ersten Oktanten. Die Punkte $K(0 | 6 | 3)$ und $L(1 | 0 | 6)$ liegen auf Kanten. $K$ liegt in einer Ebene $T$, die parallel zu $S\colon 3x + 2y + 3z = 0$ ist. Liegt auch $L$ in $T$? \hfill (Abitur 2019 GK)
\teil Bei einem Körper liegen die Ecken $A$, $B$, $C$ in der Ebene $L_1\colon 2x + y + 2z = 8$, die Ecken $D$, $E$, $F$ in $L_2\colon 4x + 2y + 4z = 14$. Begründe, dass $L_1$ und $L_2$ keine gemeinsamen Punkte haben. \hfill (Abitur 2023 GK)
\teil Eine Halle hat die quadratische Grundfläche $ABCD$ mit 6 m Seitenlänge in der $xy$-Ebene, $A$ im Ursprung, $B$ auf der $x$-Achse, $D$ auf der $y$-Achse (1 LE = 1 m). Ein Zuschauerboden liegt in der Ebene $E_B$, die parallel zu $P\colon x - y + 6z = 30$ durch $B$ verläuft. Gib $E_B$ an und bestimme die größte Höhe des Bodens über der Grundfläche. \hfill (Abitur 2020 GK)
\teil Ein dreiseitiges Prisma hat die Grundfläche $ABC$ in der Ebene $L\colon x + 2y + 2z = 3$ und die Deckfläche $DEF$; die Seitenkante $BE$ verbindet $B(1 | 1 | 0)$ mit $E(4 | 4 | 6)$. Eine zu $L$ parallele Ebene $M$ teilt das Prisma so, dass der Teilkörper, der $E$ enthält, doppelt so groß ist wie der andere. Bestimme eine Gleichung von $M$. \hfill (Abitur 2020 GK)
\end{teile}
\end{aufgabe}

%% TODO Titel: Ich-kann-Satz fehlt in der Bank (Platzhalter: Kettenname) – Nr. 18
%% TODO Anweisung: ein Satz über der ersten Teilaufgabe fehlt in der Bank – Nr. 18
\begin{aufgabe}{Parallele Ebenen über die Normalenvektoren auswählen und über die Konstante der Koordinatengleichung der Höhe nach zuordnen}
\begin{teile}
\teil Ein Turm hat drei Dachetagen; die Dachflächen der Südseite sind parallel ($z$ zeigt nach oben, 1 LE = 1 m). Die untere liegt in $E\colon 2x + 5z = 40$, die mittlere und die obere in zwei der Ebenen I: $2x + 4z = 40$, II: $2x + 5z = 30$, III: $2x + 5z = 50$, IV: $4x + 5z = 40$, V: $2x + 5z = 60$. Ordne zu und begründe. \hfill (Abitur 2017 GK)
\end{teile}
\end{aufgabe}

%% TODO Titel: Ich-kann-Satz fehlt in der Bank (Platzhalter: Kettenname) – Nr. 19
%% TODO Anweisung: ein Satz über der ersten Teilaufgabe fehlt in der Bank – Nr. 19
\begin{aufgabe}{Parallele Ebenen – Fehler finden, Begründen, Anwendung}
\begin{teile}
\teil $T$ ist parallel zu $E\colon 2x + y - z = 3$ und geht durch $K(1 | 2 | 0)$. Paul prüft, ob $L(2 | 1 | 1)$ in $T$ liegt: „$4 + 1 - 1 = 4 \ne 3$, also nicht.“ Finde den Fehler.
\teil Begründe, warum zwei parallele Ebenen denselben Normalenvektor haben können.
\teil Ein Dach liegt in $D\colon x + 3z = 30$ (1 LE = 1 m, $z$ zeigt nach oben). Eine Dämmschicht wird parallel zum Dach durch den Punkt $(6 | 0 | 7)$ angebracht. Wie weit liegt sie über $x = 0$ senkrecht unter dem Dach?
\end{teile}
\end{aufgabe}
